\section{Proofs of the verification and transport results}\label{app:bridge}
\begin{proof}[Proof of Theorem~\ref{thm:bridge}]
For any admissible control $a$, stopped It\^o's formula gives
\begin{align*}
 J(t,z;a)-v(t,z)
 =\E_{t,z}\left[\int_t^\tau e^{-\rho(s-t)}\cR_{a_s}v(s,Z_s)\,ds
       +e^{-\rho(\tau-t)}(g-v)(\tau,Z_\tau)\right].
\end{align*}
One may first stop the local martingale at a bounded localization sequence. Bounded derivatives and coefficients on the stopped compact domain give integrability, and dominated convergence then yields the displayed identity. For the candidate $\pi$, the absolute value of the right side is at most
$b+\int_t^T e^{-\rho(s-t)}e(s)ds$.
For an arbitrary original control, the residual is at most $e(s)+q(s)$, because its action is covered by the continuous maximum in \eqref{eq:checks}. Hence
\[
 V(t,z)-v(t,z)\leq b+\int_t^T e^{-\rho(s-t)}\{e(s)+q(s)\}ds.
\]
Adding the lower-policy error proves \eqref{eq:bridgebound}. Nonnegativity follows from admissibility of $\pi$. If an implemented action belongs to an enclosed neighborhood, the same residual and action-gap checks hold for it by interval inclusion, so the proof is unchanged. No statement about the convergence of the process that generated $v$ or $\pi$ is used.
\end{proof}

\begin{proof}[Proof of Proposition~\ref{prop:transport}]
In the perturbed model, the candidate residual is at least $-e-d$, while the residual of any original admissible action is at most $e+q+d$. Applying the stopped identity separately in that model bounds the lower-policy and upper-value errors by the corresponding integrals. Their sum is $2e+q+2d$, not an extra improvement-gap penalty on top of those errors. Nonnegative integrands permit replacement of the perturbed discount by $\underline\rho$. The stopping-trace error is at most $b+b_\mathrm{par}$ if the settlement differs by at most $b_\mathrm{par}$. This argument neither equates the two stopping times nor assumes a pathwise coupling of different diffusions.
\end{proof}

\section{Proof of the joint state--price construction}\label{app:stateprice}
\begin{lemma}[Concavity of the covered policy functional]\label{lem:concavity}
For each stored policy satisfying \eqref{eq:budgetshift}, there is a concave function $F_j(z;k)$ of $z\in K$, affine in $k$, such that $J_k(0,z;\pi_{j,x})\geq F_j(z;k)$. Its values at the four corners have the independently computed lower endpoints used in the state--price calculation.
\end{lemma}
\begin{proof}
Put $r=u-1$ and $L=-\log c$. The determinant expression in Section~\ref{sec:state} has the sign of
\[
 r^2L^2-2r(r+1)L+2(r+1)
   =r^2L^2+2(r+1)(1-rL)>0,
\]
while $f_{cc}<0$. Therefore $D^2 f$ is negative definite over the stated region. In the covered Gaussian integral, the preference input is affine in initial preference and consumption is affine in initial wealth. Their image remains in the region just verified. Integrating this concave composition with positive Gaussian and time weights preserves concavity.

The full-horizon preference settlement is the concave quadratic
$-.02\{(u-2+\int\theta)^2+.0025\}$. By the budget shift the terminal wealth, and hence its logarithmic settlement, is independent of initial wealth. The deterministic expenditure is independent of the initial state. Subtract a uniform normal-tail and stopping allowance from this covered integral plus its terminal settlement and expenditure. The supplement proves this allowance for all $z\in K$ and $k\in[.5,8]$. The resulting exact function is $F_j$, which retains concavity and affine dependence on price. Its coefficient of $-k$ is
\[
 B_j^{\mathrm{full}}=\tfrac12\int_0^1e^{-\rho s}\theta_{j,s}^2ds.
\]
At each corner the interval quadrature encloses the exact covered integral and terminal term. Its lower endpoint after the uniform corrections is no larger than $F_j$ there. The finite quadrature rule is not assumed to reproduce the exact functional, and sampled corner values are not used without their enclosure errors.
\end{proof}

\begin{proof}[Proof of Proposition~\ref{prop:stateprice}]
Let the four corners of $K$ be $z_v$, and write $z=\sum_v\lambda_v(z)z_v$ using the nonnegative bilinear corner weights, which sum to one. Convexity of the verified upper object and concavity of the lower functional imply
\[
 V_{k_i}(0,z)\leq U_i(z)\leq\sum_v\lambda_v(z)U_i(z_v),\qquad
 J_k(0,z;\pi_{j,x})\geq F_j(z;k)\geq\sum_v\lambda_v(z)F_j(z_v;k).
\]
For $k$ between $k_i$ and $k_{i+1}$, fixed-policy affinity in price and a common admissible policy set give convexity of $V_k$. Applying the value chord at the two endpoints yields the chord between the two verified upper corner combinations.

Suppose a policy lower corner endpoint is $L_{jv}$ at $k_j$ and its expenditure belongs to $[B_j^-,B_j^+]$. Its sign-aware lower line is
\[
 L_{jv}-(k-k_j)B_j^+\quad(k\geq k_j),\qquad
 L_{jv}-(k-k_j)B_j^-\quad(k\leq k_j).
\]
Each price cell between consecutive nodes has a fixed sign for every policy. The upper chord minus this lower line is therefore an affine function $g_{jv}(k)$. For every state, the regret of policy $j$ is at most
$\sum_v\lambda_v(z)g_{jv}(k)\leq\max_v g_{jv}(k)$.
Selecting a policy minimizing this last bound proves the uniform state guarantee
\[
 R(k)=\min_j\max_v g_{jv}(k).
\]
Take all pairwise intersections of the finitely many affine functions and the price-cell endpoints. Between successive such points, their ordering does not change. The maximum for each policy and then the minimum across policies select a single affine function on that interval. Its maximum occurs at an endpoint. Exact rational evaluation at this finite set thus bounds the supremum over all real prices. A corner line is discarded only when another corner line from the same policy dominates it at both price-cell endpoints, which proves domination on the entire cell. The program records those reductions and checks the remaining intersections without a sampled-grid extrapolation.
\end{proof}

\begin{proof}[Proof of Proposition~\ref{prop:executed}]
The independent primal and dual constructions in the supplement supply, respectively, valid original-policy lower endpoints and upper functions \eqref{eq:stateupper}. Exact action and terminal-budget checks hold for every central policy and every shifted policy at the extremes of the initial-state rectangle. Affinity of the shift extends these feasibility checks to its interior. Lemma~\ref{lem:concavity} supplies the concave lower functional with a common tail and stopping allowance. All quantities in Proposition~\ref{prop:stateprice} are therefore present.

The execution uses the 18 prices
\[
 .5,\ .875,\ 1.25,\ 1.625,\ 2,\ 2.5,\ 3,\ 3.5,\ 4,\ 4.25,
 \ 4.5,\ 5,\ 5.5,\ 6,\ 6.5,\ 7,\ 7.5,\ 8.
\]
The stored state--price record checks 211 candidate breakpoints and returns the exact fraction displayed following \eqref{eq:executedprice}; upward decimal conversion gives \eqref{eq:executedstate}. The central-state calculation uses the exact stopped-policy expenditure intervals and checks 187 candidate breakpoints with a separate all-intersections implementation. Its exact worst-case fraction is
\[
 \frac{1600901779824582944359922627717}{160587876576305697877131733762048},
\]
whose upward display is \pricebound. These are computer-assisted conclusions conditional on the explicitly stated arithmetic model and analytic enclosure lemmas, not consequences of agreement between two floating-point approximations.
\end{proof}

\section{Economic resolution and the common-endpoint stress bound}
For $p<q$, every policy loses $(q-p)B(\pi)\geq0$, and $B(\pi)\leq .02\C(1)$. Taking suprema preserves monotonicity and gives
$0\leq V_p-V_q\leq .02\C(1)(q-p)$.
Combining this fact with the independently verified intervals at $p$ and $q$ gives \eqref{eq:welfare}. On any segment where the upper and lower envelope pieces are affine, both strict separation and the conservative relative-width inequality reduce to affine inequalities in $q$. Their roots, included with the exact envelope breakpoints, produce the resolution regions without a grid test.

Finally, if every nodal upper endpoint is enlarged by $\delta$ and every nodal lower endpoint is decreased by $\delta$, every upper chord is enlarged by $\delta$ and every lower policy line is decreased by $\delta$. Their maximum/minimum envelope gap increases by at most $2\delta$. This proves the recorded endpoint stress test and also its scope: it tests sensitivity to valid endpoint inflation, not the validity of upstream node proofs.
